\input{../header.tex}

\usepackage{listings}
\usepackage{array}
\newcolumntype{C}[1]{>{\centering\let\newline\\\arraybackslash\hspace{0pt}}m{#1}}
\lstset{upquote=true,basicstyle=\fontencoding{T1}\selectfont\ttfamily}
\hypersetup{
    colorlinks=true,
    linkcolor=blue,
    filecolor=magenta,      
    urlcolor=blue,
    pdfpagemode=FullScreen,
    }

\begin{document}

% Change the following values to true to show the solutions or/and the hints
\ShowSolutiontrue
\ShowConseiltrue
\ShowWarningtrue
\titre
\cours{Fonctions, mémoire et exceptions}

Le but de cette séance est d'approfondir vos connaissances en programmation. Au terme de cette séance, l'étudiant sera capable de: 
\begin{itemize}
    \item définir une fonction et l'utiliser dans un programme,
    \item utiliser des librairies contenant des fonctions prédéfinies,
    \item connaître quelle est la portée d'une variable,
    \item comprendre comment fonctionne une pile d'exécution (call stack),
    \item gérer des exceptions.
\end{itemize}


\section{Variables et Fonctions}
\begin{Exercice}[5 minutes] \textbf{Les fonctions (fonctions basiques) (Java \textbf{et} Python)}\\
    Définissez une fonction nommée \lstinline{ping()} qui, lorsqu'elle est appelée, affiche ``pong''. \\
     
      \begin{conseil}
        Référez vous aux diapositives du cours pour la création et l'appel des fonctions.    
      \end{conseil}
      \begin{solution}
      
          \textbf{Python}:
            \lstinputlisting{ressources/Solutions/Question1.py}
          
          \textbf{\\Java}:
            \lstinputlisting{ressources/Solutions/Question1.java}
             
      \end{solution}   
  \end{Exercice}
  
    \begin{Exercice}[5 minutes] \textbf{Les Fonctions (Fonction multiplication) (Java \textbf{et} Python)}\\
    Définissez une fonction nommée \lstinline{multiplicateur()} qui prend deux arguments \textit{multiple1} et \textit{multiple2}, les multiplie et retourne le résultat. Stockez le résultat de \lstinline{multiplicateur(2,3)} dans une variable \textit{resultat} et affichez la.   \\

    \begin{conseil}
        \begin{itemize}
            \item Référez vous aux diapositives du cours pour la création et l'appel des fonctions.
            \item Pour retourner une valeur au lieu de l'afficher, utilisez le mot-clé \lstinline{return} (pour Python et Java).
        \end{itemize}        
    \end{conseil}
    \begin{solution}
        \textbf{Python}:
        \lstinputlisting{ressources/Solutions/Question2.py}
        
        \textbf{\\Java}:
        \lstinputlisting{ressources/Solutions/Question2.java}
    \end{solution}   
\end{Exercice} 
  
  \begin{Exercice}[5 minutes] \textbf{Les Fonctions (Fonctions Aire et Périmètre) (Java \textbf{et} Python)}\\
      Définissez deux fonctions nommées \lstinline{aire()} et \lstinline{perimètre()} qui prennent un argument (\lstinline{rayon}) et renvoient respectivement l'aire et le périmètre d'un cercle. Stockez les résultats dans des variables \lstinline{aire} et \lstinline{perimetre} et affichez le contenu de ces variables.   \\
     
      \begin{conseil}
          \begin{itemize}
              \item Référez vous aux diapositives du cours pour la création et l'appel des fonctions.
              \item Pour retourner une valeur au lieu de l'imprimer, utilisez le mot clé \lstinline{return} (pour Python et Java).
              \item Pour rappel, le périmètre d'un cercle s'obtient en utilisant la formule $P = 2*\pi*r$ et l'aire s'obtient en utilisant la formule $A = r^2*\pi$.
          \end{itemize}        
      \end{conseil}
      \begin{solution}
          \textbf{Python}:
          \lstinputlisting{ressources/Solutions/Question3.py}
          
          \textbf{\\Java}:
          \lstinputlisting{ressources/Solutions/Question3.java}
      \end{solution}   
  \end{Exercice} 

  \begin{Exercice}[5 minutes] \textbf{Portée des variables} (Python)\\
    Qu'affiche le programme suivant? \\

    \lstinputlisting{ressources/Questions/Question4.py}

    
     \begin{conseil}
        \begin{itemize}
            \item Le mot-clé \lstinline{global} permet d'accéder aux variables globales (définies à l'extérieur de la fonction).
            \item Soyez attentif au type des éléments retournés par chacune des fonctions.
        \end{itemize}
     \end{conseil}
     \begin{solution}
        \lstinline{x=2}\\ 
        \lstinline{x=3}\\
        \lstinline[language={bash}]{None}\\
    \end{solution}
 \end{Exercice}

 \begin{Exercice}[10 minutes] \textbf{Portée des variables} (Java) \\
    Qu'affiche le programme suivant?
    \lstinputlisting{ressources/Questions/Question5.java}
    Que se passerait-il si on décommente la ligne 17 (enlever les //)?

    \begin{conseil}
        Tenir compte de l’endroit où chaque variable est définie et quelles opérations sont réalisées sur celle-ci.
    \end{conseil}

    \begin{solution}
        \begin{itemize}
            \item  
            \lstinline{a = 2} \\
            \lstinline{a = 2} \\
            \lstinline{b = 3} \\ 

            \item On obtiendra une erreur de compilation du programme car la variable \lstinline{b} est créée à l’intérieur de la méthode \lstinline{g()} et n’existe donc qu’au sein de celle-ci. La variable \lstinline{b} n’étant pas définie en dehors de la méthode \lstinline{g()}, celle-ci ne peut pas faire l’objet d’un appel en dehors de la fonction \lstinline{g()}.
        \end{itemize}
        \end{solution}
 \end{Exercice}


\section{Gestion des exceptions}

\begin{Exercice}[5 minutes] Qu'affiche le programme suivant? \textbf{Types d'erreurs (Python)}\\
    \lstinputlisting{ressources/Questions/Question6.py}  
     \begin{conseil}
        Utilisée sur une chaîne de caractères, la fonction \lstinline{len()} renvoie le nombre de caractères de la chaîne. 
    \end{conseil}
     \begin{solution}
        \lstinline{Le résultat de la division est: 2.5}\\
        \lstinline{On ne peut pas diviser une chaîne de caractères.}
     \end{solution}   
 \end{Exercice}

\begin{Exercice}[5 minutes] Qu'affiche le programme suivant? \textbf{Types d'erreurs (Python)}\\
    \lstinputlisting{ressources/Questions/Question7.py}
    
     \begin{conseil}
         La chaîne de caractères vide est représentée par ``'' 
    \end{conseil}
     \begin{solution}
        Nous ne pouvons pas diviser un nombre par 0.
     \end{solution}   
 \end{Exercice}

 \begin{Exercice}[5 minutes] Qu'affiche le programme suivant? \textbf{Types d'erreurs (Python)}\\
    \lstinputlisting{ressources/Questions/Question8.py}
    
     \begin{conseil}
        \begin{itemize}
            \item La fonction \lstinline{float()} permet de convertir une variable en nombre décimal.
            \item La fonction \lstinline{int()} permet de convertir une variable en nombre entier.
        \end{itemize}  
    \end{conseil}

     \begin{solution}
        Nous ne pouvons pas convertir une chaîne de caractères.
     \end{solution}

\end{Exercice}
 
\begin{Exercice}[5 minutes] \textbf{Gestion d'erreurs (Java)} 
    Le programme suivant est incorrect. Que devez-vous modifier pour qu'il fonctionne correctement?
    \\
    
    \lstinputlisting{ressources/Questions/Question9.java}

    
     \begin{conseil}
        Référez-vous à la diapositive 24 du cours de cette semaine.           
     \end{conseil}
     \begin{solution}
        À l'intérieur de la fonction \lstinline{division}, lorsque b est égal à 0, le programme lève une exception de type \lstinline{ArithmeticException}. Mais dans le \lstinline{main}, l'erreur qui est interceptée est \lstinline{IndexOutOfBoundsException}. Pour corriger cette erreur, une des solutions est de remplacer \lstinline{IndexOutOfBoundsException} par \lstinline{ArithmeticException}.
     \end{solution}   
 \end{Exercice}

\begin{Exercice}[5 minutes] \textbf{Gestion d'erreurs (Java)} 
    \begin{enumerate}
    \item Qu'affiche le programme suivant supposant que l'utilisateur saisit l'entrée suivante? \\

    \begin{itemize}
        \item Entrez le premier entier: 12
        \item Entrez le deuxième entier: 2
        \item Entrez l'opération (+, -, *, /): *
    \end{itemize}

    \begin{conseil}
            Consultez \href{https://moodle.unil.ch/pluginfile.php/3705947/mod_resource/content/8/jswitch.pdf}{ce document} sur le \texttt{switch} en Java. Référez-vous à la diapositive 23 du cours. 
     \end{conseil}

     \item Que devez-vous modifier pour qu'il fonctionne?
    \end{enumerate}
    \lstinputlisting{ressources/Questions/Calculator.java}
    \begin{solution}
    \texttt{Entrez le premier entier: 12 \\ 
Entrez le deuxième entier: 2 \\
Entrez l'opération (+, -, *, /): * \\ 
24.0\\
6.0\\
Exception in thread "main" java.lang.IllegalArgumentException: L'opération n'est pas valide \\
        at Calculator.main(Calculator.java:28)}
    \end{solution}
\begin{solution}
Il faut simplement ajouter un \texttt{break;} après l'affichage du résultat d'une opération. 
\end{solution}

 \end{Exercice}

 \begin{Exercice}
   Qu'affiche le programme (question11.py) suivant supposant que l'utilisateur exécute le programme de la manière suivante dans le terminal? \\ \newline 
   \texttt{python3 question11.py 9} \\ 
   \begin{conseil}
   Référez-vous à la diapositive 10 du cours.
   \end{conseil}
   \lstinputlisting{ressources/Questions/question11.py}
   \begin{solution}
   
\texttt{Une erreur est survenue \\
Valeur finale de balance: 20
    }
   \end{solution}
 \end{Exercice}

\end{document}
